% © 2026 Ing. David Jaroš — CC BY-NC-ND 4.0
%
% This work is licensed under a Creative Commons Attribution-NonCommercial-NoDerivatives
% 4.0 International License (CC BY-NC-ND 4.0).
%
% License History: Earlier drafts (up to v0.3) were released under CC BY 4.0.
% From v0.4 onward, all material is released under CC BY-NC-ND 4.0 to protect
% the integrity of the theoretical work during ongoing academic development.
%
% See LICENSE.md for full license text.
%
% File: canonical/alpha/prime_factorization_entropy_potential.tex
% Purpose: Mathematically precise interpretation of the alpha effective potential
%          in terms of prime-factorization information entropy.
%          Provides candidate forms for the entropic term, compares their
%          stationary conditions, and clearly labels each claim.
% Date: 2026-05-11
% Related: canonical/alpha/nlogn_origin_analysis.tex
%          canonical/alpha/veff_corrected.tex
%          canonical/alpha/ALPHA_MASTER_STATUS.md
%          tools/compare_alpha_potentials.py

% UBT-AI-PROVENANCE-BEGIN
% schema: ubt-ai-provenance/v1
% tier: B_machine_verified
% ai_assistance: disclosed
% human_review: machine-verification
% editorial_responsibility: Ing. David Jaroš
% policy: ../../AI_PROVENANCE.md
% notice: Machine-verified against named sources or verifiers; individual attestation is not claimed.
% UBT-AI-PROVENANCE-END

\ifdefined\INCLUDEMODE
\else
  \documentclass[12pt,a4paper]{article}
  \usepackage[a4paper, margin=2.5cm]{geometry}
  \usepackage{amsmath, amssymb, amsthm}
  \usepackage{mathtools}
  \usepackage{hyperref}
  \usepackage{xcolor}
  \usepackage{booktabs}
  \usepackage{mdframed}
  \usepackage{tcolorbox}
  \tcbuselibrary{skins,breakable}
  \usepackage{array}

  \newtheorem{theorem}{Theorem}[section]
  \newtheorem{lemma}[theorem]{Lemma}
  \newtheorem{proposition}[theorem]{Proposition}
  \newtheorem{corollary}[theorem]{Corollary}
  \theoremstyle{definition}
  \newtheorem{definition}[theorem]{Definition}
  \theoremstyle{remark}
  \newtheorem{remark}[theorem]{Remark}

  \newmdenv[backgroundcolor=yellow!20, linecolor=orange,          linewidth=1.5pt]{gapbox}
  \newmdenv[backgroundcolor=green!15,  linecolor=green!60!black,  linewidth=1.5pt]{resultbox}
  \newmdenv[backgroundcolor=red!8,     linecolor=red!60,          linewidth=1.5pt]{warnbox}
  \newmdenv[backgroundcolor=blue!6,    linecolor=blue!40,         linewidth=1.5pt]{infobox}
  \newmdenv[backgroundcolor=gray!8,    linecolor=gray!60,         linewidth=1pt  ]{statusbox}

  \newcommand{\BB}{\mathbb{B}}
  \newcommand{\CC}{\mathbb{C}}
  \newcommand{\RR}{\mathbb{R}}
  \newcommand{\HH}{\mathbb{H}}
  \newcommand{\ZZ}{\mathbb{Z}}
  \newcommand{\NN}{\mathbb{N}}
  \newcommand{\Neff}{N_{\mathrm{eff}}}
  \newcommand{\Veff}{V_{\mathrm{eff}}}
  \newcommand{\BRam}{B_{\mathrm{Ram}}}

  % Status labels
  \newcommand{\STDLZ}{\textcolor{gray!70!black}{\textbf{[STD/L0]}}}
  \newcommand{\INTERP}{\textcolor{blue!60!black}{\textbf{[INTERP]}}}
  \newcommand{\DERIVCAND}{\textcolor{orange!80!black}{\textbf{[DERIVATION CANDIDATE]}}}
  \newcommand{\OPEN}{\textcolor{orange}{\textbf{[OPEN]}}}
  \newcommand{\PROVED}{\textcolor{green!60!black}{\textbf{[PROVED]}}}
  \newcommand{\COND}{\textcolor{blue}{\textbf{[COND]}}}

  \title{\textbf{Prime-Factorization Information Entropy and the\\
         Alpha Effective Potential}\\
         \large Structural Interpretation of the $n\log n$ Term in $V_{\mathrm{eff}}$}
  \author{Ing.\ David Jaro\v{s}}
  \date{2026-05-11}

  \begin{document}
  \maketitle
  \tableofcontents
  \newpage
\fi

%──────────────────────────────────────────────────────────────────────────────
\section{Overview and Status Discipline}
\label{sec:overview}
%──────────────────────────────────────────────────────────────────────────────

\begin{warnbox}
\textbf{Status discipline (mandatory):}
\begin{itemize}
  \item Alpha ($\alpha$) is \textbf{NOT DERIVED}.
  \item The effective coupling $B \approx 46.28$ is \textbf{NOT DERIVED} from the UBT
        action $S[\Theta]$. Gap G137-B remains \textbf{OPEN}.
  \item $B_{\mathrm{Ram}}$ defined below is a \textbf{numerical reference value} only,
        not a first-principles derivation.
  \item The prime-factorization entropy interpretation is labeled \INTERP{} throughout.
  \item Exact number-theoretic identities are labeled \STDLZ{}.
  \item Replacement of $n\log n$ by $\log\Gamma(n+1)$ is labeled \DERIVCAND{}.
\end{itemize}
\end{warnbox}

This document provides a mathematically precise structural interpretation of the
$n\log n$ term appearing in the UBT alpha-sector effective potential
$\Veff(n) = n^2 - B\cdot n\log n$.
The interpretation connects this term to prime-factorization information entropy
via the Legendre formula for $\log(n!)$ and the Stirling expansion of $\log\Gamma(n+1)$.

\textbf{Goal}: make the potential less \emph{ad hoc} by exhibiting the entropic
structure of the coefficient $n\log n$, \emph{without} claiming that this closes
Gap G137-B or derives $\alpha$.

%──────────────────────────────────────────────────────────────────────────────
\section{Motivation: the Ad Hoc Appearance of $n\log n$}
\label{sec:motivation}
%──────────────────────────────────────────────────────────────────────────────

The current best-route potential (see \texttt{canonical/alpha/alpha\_best\_route.tex}) is
\begin{equation}
  \Veff(n) = n^2 - B \cdot n\log n,
  \label{eq:veff_standard}
\end{equation}
where $n \in \NN$ is the winding-mode number and $B > 0$ is an effective coupling.
This form appears structurally motivated at the level of the free-energy counting
argument (divisor sums, one-loop determinants), but the $n\log n$ combination
looks like a product of a linear term and a logarithm without further justification.

A more structurally transparent form involves the \emph{log-factorial} or, for
continuous $n$, the log-Gamma function.  The key observation is:
\begin{equation}
  n\log n \approx \log(n!) + n - \tfrac{1}{2}\log(2\pi n) + O(1/n),
\end{equation}
by Stirling's formula.  Therefore, up to sub-leading corrections, $n\log n$ is
the \emph{leading Stirling term of $\log\Gamma(n+1)$}, which in turn admits
an exact prime-factorization (Legendre) decomposition.

%──────────────────────────────────────────────────────────────────────────────
\section{Exact Number-Theoretic Identities}
\label{sec:identities}
%──────────────────────────────────────────────────────────────────────────────

\subsection{Log-factorial and log-Gamma}
\label{sec:logfac}

\begin{resultbox}
\textbf{Identity 1} \STDLZ{}.  For every positive integer $n$,
\begin{equation}
  \log\Gamma(n+1) = \log(n!).
  \label{eq:logGamma_factorial}
\end{equation}
\end{resultbox}

This is the standard definition of the Gamma function: $\Gamma(n+1) = n!$
for $n \in \NN_0$.

\subsection{Legendre's formula (prime-factorization decomposition)}
\label{sec:legendre}

\begin{resultbox}
\textbf{Identity 2 — Legendre's formula} \STDLZ{}.  For any positive integer $n$,
\begin{equation}
  \log(n!) = \sum_{\substack{p \text{ prime} \\ p^m \leq n \\ m \geq 1}}
             \left\lfloor \frac{n}{p^m} \right\rfloor \log p,
  \label{eq:legendre}
\end{equation}
where the sum runs over all prime powers $p^m \leq n$.
\end{resultbox}

\begin{proof}
The exact power of a prime $p$ in $n!$ is $\sum_{m=1}^{\infty}\lfloor n/p^m \rfloor$
(Legendre 1808; only finitely many terms are non-zero).  Summing $\log p$ for each
prime $p$ yields~\eqref{eq:legendre}.
\end{proof}

\textbf{Remark}: Equation~\eqref{eq:legendre} is a standard result in number theory
(\STDLZ{}) and does not require any input from UBT.

\subsection{Stirling expansion}
\label{sec:stirling}

\begin{resultbox}
\textbf{Identity 3 — Stirling expansion} \STDLZ{}.
\begin{equation}
  \log\Gamma(n+1) = n\log n - n + \tfrac{1}{2}\log(2\pi n) + O(1/n),
  \quad n \to \infty.
  \label{eq:stirling}
\end{equation}
\end{resultbox}

Equivalently, the combination $\log\Gamma(n+1) + n$ has the expansion
\begin{equation}
  \log\Gamma(n+1) + n = n\log n + \tfrac{1}{2}\log(2\pi n) + O(1/n).
  \label{eq:logGammaPlusn}
\end{equation}
The \emph{leading term} is $n\log n$ in both cases.

\subsection{Digamma (derivative of log-Gamma)}
\label{sec:digamma}

\begin{resultbox}
\textbf{Identity 4} \STDLZ{}.  The derivative of $\log\Gamma(n+1)$ with respect to
(continuous) $n$ is the digamma function $\psi$:
\begin{equation}
  \frac{d}{dn}\log\Gamma(n+1) = \psi(n+1),
  \label{eq:digamma_def}
\end{equation}
with asymptotic expansion
\begin{equation}
  \psi(n+1) = \log n + \frac{1}{2n} - \frac{1}{12n^2} + O(n^{-4}),
  \quad n \to \infty.
  \label{eq:digamma_asymp}
\end{equation}
\end{resultbox}

%──────────────────────────────────────────────────────────────────────────────
\section{Information-Theoretic Interpretation}
\label{sec:info_theory}
%──────────────────────────────────────────────────────────────────────────────

\begin{infobox}
\textbf{Interpretation} \INTERP{}.  Consider the \emph{prime alphabet}
$\mathcal{P} = \{2, 3, 5, 7, 11, \ldots\}$ with symbol weights $w_p = \log p$.
Each positive integer $k \leq n$ can be encoded as a \emph{composite symbol}
over this alphabet via its prime factorization $k = \prod_p p^{a_p(k)}$.

Define the \emph{logarithmic information content} of $k$ as
$H(k) = \log k = \sum_p a_p(k)\log p$.

The total information content of all integers $1, 2, \ldots, n$ is
\begin{equation}
  S_{\log}(n)
  := \sum_{k=1}^{n} H(k)
   = \sum_{k=1}^{n} \log k
   = \log(n!).
\end{equation}

By Identity 2 (Legendre), this equals
\begin{equation}
  S_{\log}(n) = \sum_{\substack{p \text{ prime},\, m \geq 1 \\ p^m \leq n}}
                \left\lfloor \frac{n}{p^m} \right\rfloor \log p,
  \label{eq:Slog_legendre}
\end{equation}
i.e.\ the sum of $\log p$ over all prime-power occurrences in the factorizations
of integers $1$ through $n$.

\emph{Therefore}: the entropic term $\log(n!) = \log\Gamma(n+1)$ is the total
logarithmic information content of the integers $1, \ldots, n$ when encoded
over the weighted prime alphabet $(\mathcal{P},\log)$.  Its leading behaviour
is $n\log n$ (Stirling).
\end{infobox}

This interpretation is entirely \INTERP{}: it does not derive $B$ from $S[\Theta]$,
nor does it provide a UBT mechanism for the entropic term.  It explains \emph{why}
$n\log n$ is a natural quantity (it is the leading term of a combinatorial/information
entropy), but the connection to UBT dynamics remains Gap G137-B.

%──────────────────────────────────────────────────────────────────────────────
\section{Three Candidate Potential Forms}
\label{sec:candidates}
%──────────────────────────────────────────────────────────────────────────────

We compare three forms of the effective potential and their stationary conditions.

%──────────────────────────────────────────────────────────────────────────────
\subsection{Form V1: Original potential}
\label{sec:V1}
%──────────────────────────────────────────────────────────────────────────────

\begin{equation}
  V_1(n) = n^2 - B\cdot n\log n.
  \label{eq:V1}
\end{equation}

\textbf{Derivative}:
\begin{equation}
  V_1'(n) = 2n - B(\log n + 1).
  \label{eq:V1_deriv}
\end{equation}

\textbf{Stationary condition} $V_1'(n^*) = 0$:
\begin{equation}
  B = \frac{2n^*}{\log n^* + 1}.
  \label{eq:V1_stat}
\end{equation}

For $n^* = 137$: $B = 2\times137/(\log 137 + 1) \approx 274/5.920 \approx 46.28$.

%──────────────────────────────────────────────────────────────────────────────
\subsection{Form V2: Gamma-refined potential}
\label{sec:V2}
%──────────────────────────────────────────────────────────────────────────────

\DERIVCAND{}

\begin{equation}
  V_2(n) = n^2 - B\cdot\log\Gamma(n+1).
  \label{eq:V2}
\end{equation}

\textbf{Motivation}: Replace $n\log n$ by the exact function $\log\Gamma(n+1)$,
whose prime-factorization decomposition is given by Legendre's formula.

\textbf{Derivative} (using Identity 4):
\begin{equation}
  V_2'(n) = 2n - B\cdot\psi(n+1).
  \label{eq:V2_deriv}
\end{equation}

\textbf{Stationary condition} $V_2'(n^*) = 0$:
\begin{equation}
  B = \frac{2n^*}{\psi(n^*+1)}.
  \label{eq:V2_stat}
\end{equation}

Asymptotically, $\psi(n+1) \approx \log n + O(1/n)$ (Identity 4), so
\begin{equation}
  B \approx \frac{2n^*}{\log n^*}.
  \label{eq:V2_stat_asymp}
\end{equation}

\begin{warnbox}
\textbf{Caution}: The stationary condition of $V_2$ differs from that of $V_1$.
$V_1'(n) = 2n - B(\log n + 1)$, whereas $V_2'(n) = 2n - B\cdot\psi(n+1)
\approx 2n - B\log n$.  This shifts the effective $B$ by a factor of approximately
$(\log n)/(\log n + 1)$.  For $n^* = 137$: $\log 137 \approx 4.920$, so the ratio
is $\approx 0.831$ — a $17\%$ shift.  \emph{Replacing $V_1$ by $V_2$ silently
changes the numerical prediction of $B$ for a given $n^*$.}
\end{warnbox}

For $n^* = 137$: $B = 2\times137/\psi(138) \approx 274/4.924 \approx 55.64$.
This differs significantly from the phenomenological $B_{\mathrm{phenom}} \approx 46.28$,
making $V_2$ alone a \emph{less accurate} candidate if we require $n^* = 137$.

%──────────────────────────────────────────────────────────────────────────────
\subsection{Form V3: Recommended refined form}
\label{sec:V3}
%──────────────────────────────────────────────────────────────────────────────

\DERIVCAND{}

\begin{equation}
  V_3(n) = n^2 - B\cdot\bigl(\log\Gamma(n+1) + n\bigr).
  \label{eq:V3}
\end{equation}

\textbf{Motivation}: By Identity 3 (equation~\eqref{eq:logGammaPlusn}),
\begin{equation}
  \log\Gamma(n+1) + n = n\log n + \tfrac{1}{2}\log(2\pi n) + O(1/n).
\end{equation}
This preserves the leading $n\log n$ term while embedding it in the prime-factorization
entropy structure.

\textbf{Derivative}:
\begin{equation}
  V_3'(n) = 2n - B\bigl(\psi(n+1) + 1\bigr).
  \label{eq:V3_deriv}
\end{equation}

\textbf{Stationary condition} $V_3'(n^*) = 0$:
\begin{equation}
  B = \frac{2n^*}{\psi(n^*+1) + 1}.
  \label{eq:V3_stat}
\end{equation}

Asymptotically,
\begin{equation}
  B \approx \frac{2n^*}{\log n^* + 1},
  \label{eq:V3_stat_asymp}
\end{equation}
which \emph{matches the stationary condition of $V_1$} to leading order.

For $n^* = 137$:
\begin{align}
  \psi(138) &\approx 4.924, \\
  B &= 2\times137/(4.924 + 1) \approx 274/5.924 \approx 46.25,
\end{align}
which closely reproduces the phenomenological value $B_{\mathrm{phenom}} \approx 46.28$.

\begin{resultbox}
\textbf{Key result}: $V_3$ combines (i) the exact prime-factorization entropy structure
of $\log\Gamma(n+1)$ with (ii) preservation of the original stationary condition
$B \approx 2n/(\log n + 1)$ at leading order.  This makes $V_3$ the preferred
\DERIVCAND{} refinement of the original $V_1$.
\end{resultbox}

%──────────────────────────────────────────────────────────────────────────────
\section{Summary Comparison}
\label{sec:comparison}
%──────────────────────────────────────────────────────────────────────────────

\begin{center}
\begin{tabular}{@{}llll@{}}
\toprule
& $V_1 = n^2 - Bn\log n$ & $V_2 = n^2 - B\log\Gamma(n+1)$ & $V_3 = n^2 - B(\log\Gamma(n+1)+n)$ \\
\midrule
Derivative & $2n - B(\log n + 1)$ & $2n - B\psi(n+1)$ & $2n - B(\psi(n+1)+1)$ \\
Stat.\ cond.\ (exact) & $B = \frac{2n}{\log n+1}$ & $B = \frac{2n}{\psi(n+1)}$ & $B = \frac{2n}{\psi(n+1)+1}$ \\
$B$ at $n^*=137$ & $\approx 46.28$ & $\approx 55.64$ & $\approx 46.25$ \\
Entropy form & approx.\ Stirling leading & exact log-Gamma & exact log-Gamma$+n$ \\
Stat.\ cond.\ (asymptotic) & $\frac{2n}{\log n+1}$ & $\frac{2n}{\log n}$ & $\frac{2n}{\log n+1}$ \\
Match to $B_\mathrm{phenom}$ & exact by def. & $\approx -17\%$ & $< 0.1\%$ \\
Status & [L1] & \DERIVCAND{} & \DERIVCAND{} \\
\bottomrule
\end{tabular}
\end{center}

%──────────────────────────────────────────────────────────────────────────────
\section{Relation to the Alpha Problem and Gap G137-B}
\label{sec:alpha_gap}
%──────────────────────────────────────────────────────────────────────────────

\begin{gapbox}
\textbf{Gap G137-B} \OPEN{}.  The coupling $B \approx 46.28$ required for
$n^* = 137$ has not been derived from the UBT action $S[\Theta]$ without using
$\alpha$ as an input.  This gap is \textbf{not closed} by the prime-factorization
entropy interpretation.

The present document provides only:
\begin{enumerate}
  \item An exact number-theoretic decomposition of the entropic term (\STDLZ{}).
  \item An information-theoretic \emph{interpretation} of why $n\log n$ is natural
        (\INTERP{}).
  \item A candidate potential refinement $V_3$ that preserves the stationary condition
        while embedding the entropy structure (\DERIVCAND{}).
\end{enumerate}

None of the above constitutes a derivation of $B$ from $S[\Theta]$, nor a proof
that $\alpha^{-1} = 137$.  Alpha is \textbf{NOT DERIVED}.
\end{gapbox}

The prime-factorization entropy interpretation does, however, provide a structural
explanation for the $n\log n$ form: it is the leading Stirling term of the total
logarithmic information content of integers $1, \ldots, n$ under the prime alphabet.
This replaces an apparently ad hoc combination with a well-defined combinatorial
quantity, strengthening the structural motivation for the potential form without
overstating the derivation status.

%──────────────────────────────────────────────────────────────────────────────
\section{Numerical Verification}
\label{sec:numeric}
%──────────────────────────────────────────────────────────────────────────────

The script \texttt{tools/compare\_alpha\_potentials.py} computes, for each of $V_1$,
$V_2$, $V_3$:
\begin{itemize}
  \item the stationary-point equation and its numerical solution as a function of $B$;
  \item the value of $B$ required for $n^* = 137$;
  \item the relative error in the stationary condition compared to the phenomenological
        value $B_{\mathrm{phenom}} \approx 46.28$.
\end{itemize}

The reference value used numerically is
\begin{equation}
  \BRam := 12^{3/2} \cdot 2^{1/8} \cdot \vartheta_3(0\mid i)^{1/4},
  \label{eq:BRam}
\end{equation}
which appears as a numerical reference in some UBT computations.
This expression is a \textbf{numerical observation only}; it has been explicitly
rejected as a first-principles derivation of the $n\log n$ coefficient
(see \texttt{canonical/alpha/ALPHA\_MASTER\_STATUS.md}, deprecated-claim register).

%──────────────────────────────────────────────────────────────────────────────
\ifdefined\INCLUDEMODE
\else
  \end{document}
\fi
